\documentclass[10pt]{article}
\usepackage[margin=2.2cm]{geometry}
\usepackage{microtype}
\usepackage{url}
\usepackage{hyperref}
\hypersetup{hidelinks}
\setlength{\parindent}{0pt}
\setlength{\parskip}{5pt}

\title{\vspace{-1.2cm}Summary of Changes\\[2pt]
\large Distance-Preserving Digests: A Primitive for BFT Consensus (BLOCKCHAIN'26 camera-ready)}
\author{Ryan Mercier}
\date{}

\begin{document}
\maketitle
\vspace{-1.1cm}

References give the location in the submitted IEEEtran version first and in
this Springer camera-ready second; condensation to the 10-page limit turned
several IEEE subsections into run-in paragraphs inside renumbered sections.

\section*{Reviewer 1}

\textbf{1. No heading directly followed by another heading.} Every section
now opens with a substantive lead-in paragraph stating what problem the
section addresses and how its parts connect. Lead-ins were added at the
former junctions: the primitive (was Sect.~III, now Sect.~4), the consensus
protocol (was Sect.~IV, now Sect.~5), cross-shard consistency (was
Sect.~V, now Sect.~6), the evaluation (was Sect.~VII, now Sect.~7), and the
security analysis (was Sect.~VIII, now Sect.~8).

\textbf{2. Related work should follow the introduction.} Related work moved
from the back (was Sect.~IX) to directly after the introduction (now
Sect.~2). The redundancy this exposed with the introduction (LSH novelty,
Ethereum committee mechanics, cross-shard coordination detail) was merged
into the relocated section; the trims are logged in \texttt{CUTS.md} in the
repository.

\textbf{3. Add implementation details.} An implementation paragraph now
closes the evaluation (was Sect.~X, now the ``Implementation'' paragraph of
Sect.~7): digest computation cost and layout (one SHA-512 plus eight
modular reductions per transaction, vectorized sums, 64 bytes on the wire),
the two interchangeable BLS code paths (real BLS12-381 via py-ecc in the
demo, hash-based mocks preserving the 96-byte signature and 48-byte
public-key wire sizes in benchmarks) with their cross-validation
(identical per-block message totals on the small-$N$ demo under both), and
the aggregation-side optimizations in the codebase (single-push bloom-diff
sync, direct aggregation of speculative Phase~1 commitments into the
fast-path certificate, per-leaf splitting of tree aggregation). Only what
exists in the cited repository is described.

\section*{Reviewer 2}

\textbf{1. Abstract: the 8-dimensional digest space needs background.} The
abstract now explains the object before using it: each transaction hashes
to a point in an 8-dimensional space, a validator's digest is the sum of
its transactions' points, and the Euclidean distance between two digests is
therefore proportional to the number of differing transactions. The
abstract stays within 150--250 words.

\textbf{2. Motivate the static adversary model.} The introduction's closing
paragraph now states that static corruption is the standard first-analysis
setting for this protocol class, the model in which the classical safety
arguments for PBFT, HotStuff, and Tendermint are stated, and that it
isolates the primitive's contribution from adaptive corruption; adaptive
adversaries are declared out of scope with a pointer to the limitations
discussion (now Sect.~9).

\textbf{3. Sect.~III.D numbers must be derivable.} The liveness bound (now
the ``Threshold calibration and liveness bound'' paragraph of Sect.~4)
derives every constant in prose: $k_{\max}{=}2$ is the partial-observation
model the calibration samples; the sampled 99th percentile ${\approx}4.1$
times the $1.2$ margin gives $\tau \approx 4.9$; the conditional exceedance
$0.005$ is measured over the same sample; multiplying by $p_{\mathrm{miss}}
= 0.37$ gives $p = 0.00185$; Hoeffding with $t = 1/3 - p \approx 0.33$
(so $2t^2 \approx 0.22$) yields $\exp(-0.22N)$, evaluated as $\exp(-22)
\approx 3{\times}10^{-10} < 10^{-9}$ at $N{=}100$ and $\exp(-220) <
10^{-95}$ at $N{=}1000$.

\textbf{4. Overview paragraph for the protocol section.} Sect.~5 (was
Sect.~IV) opens by stating what each phase does, that the phases collapse
into one round trip when enough validators already hold the block, and why
the signature phase can be accelerated but never skipped.

\textbf{5. Overview paragraph for the evaluation section.} Sect.~7 (was
Sect.~VII) opens with the three questions the evaluation answers (message
load at scale, the BLS bottleneck, the cost to Byzantine tolerance) and the
measurement discipline (a counter on every simulated send, identical byte
constants across protocols).

\textbf{6. Informal phrasing in the cross-shard section.} In Sect.~6 (was
Sect.~V): ``Lower latency, but receipts scale with transaction volume''
became ``\ldots achieves lower latency than two-phase commit at the price
of receipt volume that grows linearly with cross-shard transaction
volume''; ``tells you identical or not-identical'' became ``hash-based
verification indicates only whether two states are identical, revealing
nothing about the magnitude of divergence.'' The section's remaining
fragments (the table lead-in, the multi-shard enumeration) were rewritten
as complete sentences.

\textbf{7. HotStuff referenced without introduction.} At first mention
(Sect.~1) HotStuff is now described self-contained: a leader-based BFT
protocol whose rotating leader drives each block through three voting
rounds, collecting each round's votes into a threshold-signature quorum
certificate proving $2N/3$ support. Later references build on this; nothing
beyond what the comparisons use was added.

\section*{Additional changes}

\textbf{(a) Template.} Reformatted from IEEEtran two-column to the Springer
template required by the congress (\texttt{llncs} class, \texttt{splncs04}
bibliography, keywords after the abstract, no page numbers), renumbering
all sections, figures, and tables.

\textbf{(b) Author-identified fast-path correction.} The submitted version
issued the fast-path certificate from Phase~1 digest agreement alone
(near-zero cluster variance), backed by no signatures, inconsistent with
the paper's own analysis (Theorem~1's premise is $2N/3$ signed commitments;
Sect.~VIII.D states Phase~1 contributes no cryptographic security). The
mechanism is repaired, not softened: Phase~1 messages now carry a BLS
commitment on $(\mathit{height}, H(B_{\mathrm{prop}}))$, attached iff the
validator's local state matches the proposal ($64{+}25{+}96$-byte payload
from matching validators), and the fast path triggers only when the
aggregator collects ${\ge}2N/3$ valid commitments on one hash in round one,
aggregating them into a certificate byte-identical to the Phase~2 object
(96-byte aggregate plus $N/8$-byte bitmap). Single-round finality is thus
covered by Theorem~1 verbatim, and an aggregator can at most withhold a
certificate (liveness), as noted after Corollary~1. On a miss, the round
falls back to full Phase~2 collection exactly as submitted; Phase~1
commitments may be reused there, a saving the evaluation conservatively
does not model. The fast path applies to the flat protocol; carrying
speculative commitments up the tree is deferred.
$(1-p_{\mathrm{miss}})^N$ is now labelled a lower bound on engagement
(matching all honest validators, where only $2N/3$ commitments are
required). The mechanism follows the speculative lineage of Zyzzyva (Kotla
et al., SOSP 2007) and SBFT (Gueta et al., DSN 2019), both now cited. The
Byzantine-sweep figure (was Fig.~9, now Fig.~2) was regenerated by
re-running the repository's figure script so its bandwidth panel reflects
the 96-byte Phase~1 commitment; a pinned-clock comparison verified that
success rates, message counts, and the other protocols' bandwidth are
bit-identical, and at $p_{\mathrm{miss}}{=}0.37$ the fast path does not
engage, so all other reported numbers are unchanged. The repository
implementation was updated to match, and the protocol-flow figure was
redrawn (now script-generated in the repository) so its decision node,
Phase~1 message contents, and fast-path terminal show the corrected
mechanism.

\textbf{(c) Condensation.} The mechanical Springer conversion ran to 17
pages before the requested additions, so the camera-ready is condensed to
the 10-page limit: secondary derivations became one-line results with
pointers to the extended version (technical report at
\url{https://github.com/RyanMercier/Proxima}), figures whose plotted values
are fully stated in the retained tables and prose were removed, the
messages and latency tables were merged, and several subsections became
run-in paragraphs. No reviewed number, table value, or plotted datum was
changed; no technical content beyond (b) was added. The full cut log is
\texttt{CUTS.md}.

\end{document}
